\chapter{Introduction}
\label{chp:intro}

This manual describes the hardware of the Wildbits \jrtwo\ and \ktwo\ running the FNX6809 core: what each device does, which registers control it, and how to program it from 6809 assembly language. Each chapter covers one part of the machine and gives working examples.

The manual doesn't teach 6809 programming. The instruction set is the one used in the Tandy Color Computer, the Dragon and many other machines, and books written for those machines teach it well. The examples are commented so you can follow them if you are still learning.

\section{The machines}

The Wildbits \jrtwo\ and \ktwo\ are the main subject of this manual.

\begin{description}
    \item[\jrtwo] a single board computer in the mini-ITX format. It needs a PS/2 keyboard, a case and a power supply of your own.
    \item[\ktwo] a complete computer with a case and a built-in keyboard. It has a few devices the \jrtwo\ lacks: a second VIA that scans the built-in keyboard, an optical keyboard interface, a WizNet network interface, a supervisor microcontroller (an RP2040) and a small logo display.
\end{description}

Both machines run the same core and share almost all their hardware, so a program written for one runs on the other unless it uses a \ktwo-only device. The text marks every difference.

\begin{note}
    The same core will be ported to the earlier \fjr\ and \fk. Most of this manual will apply to them, but they have less memory and fewer devices, so check the chapter for the device you use.
\end{note}

\section{The 6809 and the core}

Most of the machine lives in one FPGA. The core loaded into it, FNX6809, provides a cycle-accurate 6809 CPU along with VICKY, the video controller; the memory management unit (MMU); the sound chips; and most of the I/O devices. Because the CPU is part of the core, you change the CPU by loading a different core; this manual covers the 6809 core only. This edition describes core release RC24 (``Longview'').

The CPU is a 6809, not a 6309: the 6309's extra registers and instructions aren't there. It runs at 6.29~MHz. With turbo mode on, cycles that don't use the bus (the internal cycles of an instruction) run at 12.58~MHz, while every bus cycle keeps its normal timing, so programs run faster without upsetting any device. Turbo mode is on when the first boot-mode DIP switch is on; bit 0 of \verb+$FF90+ then reads 0. \verify{Check the DIP switch labels and the speed-up on both boards.}

The hardware registers are big-endian, like the 6809: a register of two or more bytes keeps its most significant byte at the lowest address. You can load and store 16-bit values with \verb+LDD+ and \verb+STD+, and a 24-bit address is a high byte followed by a 16-bit word.

\section{The parts of the machine}

The CPU addresses 64~KB through the MMU, which maps 8~KB blocks of the 2,016~KB of RAM and the device memory into eight slots. Chapter~\ref{chp:memory} explains the MMU; read it first, because every other chapter depends on it.

The rest of the hardware falls into these groups, each with its own chapter:

\begin{description}
    \item[Video] VICKY draws text, bitmaps, tiles and sprites, with a mouse pointer and line interrupts. It reads its data straight from RAM, without the CPU.
    \item[Sound] the FPGA holds SID chips, an SN76489 (the PSG) and an OPL3, mixed through the sound CODEC. \verify{A VS1053b decodes compressed audio, and a SAM2695 synthesizer plays MIDI.}
    \item[Timing] two timers, the start-of-frame and start-of-line interrupts, and a battery-backed real-time clock.
    \item[Input] PS/2 keyboard and mouse, two joystick ports through VIA~0, NES and SNES gamepads, and on the \ktwo\ the built-in keyboard.
    \item[Storage and communication] \verify{two SD card interfaces}, the serial port, the Commodore IEC serial bus, and Wi-Fi through a WizFi360 module.
    \item[Helpers] a DMA engine for copying and filling memory, integer and floating-point math units, and a random number generator.
\end{description}

Interrupts from all these devices pass through one interrupt controller (chapter~\ref{chap:interrupts}). The memory map appendix lists every register address.

\section{Raw programs and NitrOS-9}

You can program the machine in two ways, and most chapters give an example of each.

\begin{description}
    \item[Raw 6809] your program runs on the bare machine, with no operating system. It owns all the hardware: it sets up the MMU, the screen and the interrupts itself. Write this way for games, demos and tests, or to learn how the hardware works.
    \item[NitrOS-9] your program runs as a process under NitrOS-9 Level 2, a multitasking operating system with files, a shell and device drivers. The system owns the MMU and the interrupt controller, so a program asks the system to map a device block before it touches a register, and it leaves interrupts to the drivers. Most software for these machines runs under NitrOS-9.
\end{description}

Where NitrOS-9 changes how you use a device, the chapter has a section on it. Use a NitrOS-9 built for the RC24 core: older builds expect the device blocks at different addresses.

\section{The examples}

Complete versions of the examples are in the \verb+examples+ folder of the manual's repository, \url{https://github.com/jfed6000/Wildbits-F256-6809-Manual}. The folder has three parts:

\begin{description}
    \item[\texttt{examples/raw}] raw programs. Start with \verb+skeleton.asm+, which sets up the machine, clears the screen and prints a message; copy it to begin your own program. Its header explains how to build it and load it.
    \item[\texttt{examples/nos9}] NitrOS-9 programs. Start with \verb+hello.asm+. The README explains how to build them and copy them to a NitrOS-9 disk image.
    \item[\texttt{examples/include}] \verb+wildbits_regs.inc+, which defines a name for every register this manual documents. It is generated from the same data as the register tables in this manual, so the names match. NitrOS-9 programs use the names in NitrOS-9's own \verb+wildbits.d+ instead.
\end{description}

To load a raw program, connect the debug USB port to another computer and use FoenixMgr. Tell the machine to boot from RAM once, then upload the program:

\begin{verbatim}
fnxmgr --boot ram
fnxmgr --binary skeleton.bin --address 00E000
\end{verbatim}

The upload stops the CPU, writes the program into RAM and restarts the CPU, which takes its reset vector from the program. \verb+fnxmgr --boot flash+ returns to the normal boot. Chapter~\ref{chp:memory} explains what the RAM boot changes, and the debug port chapter describes the protocol FoenixMgr uses.

\section{Notation}

Notes that call out a difference between machines, or something easy to get wrong, have a bar in the margin:

\begin{note}
    Like this.
\end{note}

Numbers that start with \verb+$+ are hexadecimal, and numbers that start with \verb+%+ are binary; other numbers are decimal. A range of addresses such as \verb+$FE20+--\verb+$FE2F+ includes both ends. A bit is numbered from 0, the least significant bit, and \verb+b7+ means bit 7.

Addresses come in two kinds. A CPU address is one of the 64~KB the 6809 sees, written with four digits (\verb+$FFA0+). A physical address is a place in the machine's memory, written with six digits (\verb+$1F0000+), or as a block number and an offset (block \verb+$FC+, offset \verb+$1700+). The MMU turns one into the other.

The examples use the syntax of lwasm, the assembler in lwtools (\url{http://www.lwtools.ca/}). Anything that assembles 6809 code with Motorola syntax will need few changes. The parts of lwasm syntax the examples use are:

\begin{itemize}
    \item \verb+label equ value+ gives a name to a number.
    \item \verb+fcb+ stores bytes, \verb+fdb+ stores 16-bit words (high byte first), \verb+fcc+ stores a string and \verb+rmb+ reserves bytes.
    \item \verb+org+ sets the address of the code that follows.
    \item An expression may use \verb|+ - * /|, \verb+&+ (and), \verb+|+ (or), \verb+^+ (exclusive or) and \verb+~+ (complement). For example, \verb+anda #~INT_VKY_SOF+ clears the bit named \verb+INT_VKY_SOF+.
    \item A label that starts with \verb+@+ is local: it is known only between the two ordinary labels around it.
    \item Text after the operand, or after a \verb+*+ or \verb+;+ at the start of a line, is a comment.
\end{itemize}

The raw examples build with \verb+lwasm --6809 --format=raw+, which writes a plain binary image. The NitrOS-9 examples build with \verb+--format=os9+ and the same options as the NitrOS-9 sources; each example's Makefile has the full command.
